\section{Dispatcher}
\label{sec:dispatcher}

\subsection{Rolle und Verantwortlichkeiten}

Der Dispatcher ist die zentrale Koordinationskomponente des TaskGrid+{}-Systems.
Er nimmt Tasks von Clients entgegen, verwaltet ihren Lebenszyklus vollständig,
koordiniert die Zuweisung an Worker und stellt Ergebnisse bereit.

\begin{table}[H]
\centering
\begin{tabular}{ll}
\toprule
\textbf{Verantwortlichkeit} & \textbf{Umsetzung} \\
\midrule
Tasks empfangen und validieren             & \texttt{PostTask} (Issue \#13) \\
Tasks in Warteschlange einreihen           & \texttt{TaskQueue} (Issue \#21) \\
Worker über Namensdienst auflösen          & \texttt{NamensdienstClient} (Issue \#14) \\
Tasks an Worker dispatchen                 & \texttt{DispatchLoop} (Issue \#15) \\
Ergebnisse von Workern empfangen           & \texttt{ReturnResult} (Issue \#16) \\
Ergebnisse an Clients ausliefern           & \texttt{GetResult} (Issue \#17) \\
Timeouts erkennen und Retries steuern      & \texttt{DispatchLoop} (Issue \#18) \\
Strukturierte Logs schreiben               & \texttt{logger.py} (Issue \#19) \\
Monitoring-Status bereitstellen            & \texttt{StatusCollector} (Issue \#20) \\
Nebenläufigkeit und Race Conditions        & Locking-Strategie (Issue \#21) \\
\bottomrule
\end{tabular}
\caption{Dispatcher -- Verantwortlichkeiten nach Issue}
\end{table}

\subsection{Komponentenarchitektur}

Der Dispatcher besteht aus sechs klar getrennten Modulen, die zusammen den
vollständigen Task-Lebenszyklus abdecken:

\begin{lstlisting}[language={},caption={Dispatcher-Komponenten (src/dispatcher/)}]
src/dispatcher/
  server.py               -- Einstiegspunkt, startet alle Threads
  servicer.py             -- gRPC-Handler (PostTask, ReturnResult, GetResult, GetStatus)
  dispatch_loop.py        -- Background-Thread: Queue -> Namensdienst -> Worker
  task_store.py           -- In-Memory-Speicher fuer alle Tasks (thread-sicher)
  task_queue.py           -- FIFO-Warteschlange (queue.Queue)
  state_machine.py        -- Zustandsuebergangspruefung (8 Zustaende)
  namensdienst_client.py  -- gRPC-Client fuer Namensdienst-Abfragen
  worker_client.py        -- gRPC-Client fuer ExecuteTask an Worker
  worker_selector.py      -- Round-Robin-Selektor
  status_collector.py     -- Aggregiert Monitoring-Daten
  http_status_server.py   -- HTTP-Daemon-Thread (GET /status, Port 8080)
src/common/
  protocol.py             -- TaskState-Enum, new_task_id()
  logger.py               -- Strukturiertes Logging (log_event)
\end{lstlisting}

\subsubsection*{Startsequenz}

Beim Start (\texttt{server.py}) werden die Komponenten in folgender Reihenfolge
initialisiert und gestartet:

\begin{enumerate}
  \item \texttt{TaskStore} und \texttt{TaskQueue} anlegen (in-memory).
  \item \texttt{NamensdienstClient} verbinden (Adresse aus Env-Vars).
  \item \texttt{DispatchLoop} starten (Daemon-Thread).
  \item \texttt{StatusCollector} und \texttt{HttpStatusServer} starten (Daemon-Thread, Port 8080).
  \item gRPC-Server starten (Port 50051, ThreadPoolExecutor mit 10 Worker-Threads).
  \item SIGTERM/SIGINT-Handler registrieren (Graceful Shutdown).
\end{enumerate}

\subsection{Proto-Kompatibilität mit Elena's Worker}
\label{sec:dispatcher-proto}

Das gemeinsame Nachrichtenformat ist in \texttt{proto/taskgrid.proto}
(proto3) definiert. Da Elenas Worker-Code eigene Proto-Stubs mitbringt,
wurden alle Dispatcher-Proto-Definitionen auf Elenas Konventionen abgestimmt
(\textit{``Passe alles auf Elena an, mache Kompromisse, ändere nichts an Elenas Code''}):

\begin{table}[H]
\centering
\begin{tabular}{lll}
\toprule
\textbf{Aspect} & \textbf{Elena-Konvention} & \textbf{Dispatcher-intern} \\
\midrule
Task-ID im Proto      & \texttt{int32 task\_id}      & \texttt{str} (thread-sicherer Zähler) \\
Task-Payload-Feld     & \texttt{task\_payload}       & \texttt{task.payload} \\
Ergebnis-Methode      & \texttt{ReturnResult}        & \texttt{ReturnResult} \\
Status-Felder         & \texttt{status="COMPLETED"/"FAILED"} & \texttt{TaskState}-Enum \\
Fehlerfeld            & \texttt{error}               & \texttt{error} \\
ACK-Feld              & \texttt{Ack.success} (bool)  & \texttt{Ack.success} \\
Namensdienst-Service  & \texttt{NamingService}       & \texttt{NamingServiceStub} \\
\bottomrule
\end{tabular}
\caption{Proto-Kompatibilität Dispatcher/Elena}
\end{table}

Die Konvertierung erfolgt ausschließlich an der gRPC-Grenze:
\begin{lstlisting}[language=Python,caption={int32-Konvertierung an der gRPC-Grenze}]
# Eingehend (int32 -> str):
task_id = str(request.task_id)   # in ReturnResult, GetResult

# Ausgehend (str -> int32):
return taskgrid_pb2.TaskResponse(
    task_id=int(task_id),         # in PostTask
    status=TaskState.QUEUED.value,
)
\end{lstlisting}

\subsection{Issue \#13: POST\_TASK -- Aufgaben empfangen}
\label{sec:post-task}

\texttt{PostTask} ist die Eingangsschnittstelle für Clients.
Die Implementierung in \texttt{servicer.py} durchläuft drei Validierungsstufen,
bevor der Task akzeptiert wird:

\begin{lstlisting}[language=Python,caption={PostTask -- Validierung und Einreihung}]
# 1. Validierung
if not request.task_type:
    context.set_code(grpc.StatusCode.INVALID_ARGUMENT)
    return taskgrid_pb2.TaskResponse()

if len(request.task_type) > MAX_TYPE_LEN:      # max 32 Zeichen
    context.set_code(grpc.StatusCode.INVALID_ARGUMENT)
    return taskgrid_pb2.TaskResponse()

if len(request.task_payload) > MAX_PAYLOAD_LEN: # max 1024 Zeichen
    context.set_code(grpc.StatusCode.INVALID_ARGUMENT)
    return taskgrid_pb2.TaskResponse()

# 2. Task anlegen und speichern
task_id = new_task_id()           # thread-sicherer int32-Zaehler
task = Task(task_id=task_id, task_type=request.task_type,
            payload=request.task_payload, status=TaskState.QUEUED)
self._store.add(task)
self._queue.enqueue(task)

# 3. Antwort
return taskgrid_pb2.TaskResponse(
    task_id=int(task_id),          # int32 fuer Elena-Kompatibilitaet
    status=TaskState.QUEUED.value,
)
\end{lstlisting}

\textbf{Alle drei Fehlerfälle} sind durch Integration Tests abgedeckt und
werden im Testprotokoll dokumentiert (Abschnitt~\ref{sec:tests}).

\subsection{Issue \#14: Worker-Lookup und Round-Robin}
\label{sec:worker-lookup}

Der Dispatcher löst Worker-Adressen \textbf{ausschließlich dynamisch} auf.
Statische Adressen existieren nirgends im Code oder in Konfigurationsdateien.

\subsubsection*{Namensdienst-Client}

\texttt{NamensdienstClient} kapselt alle gRPC-Aufrufe an den Namensdienst.
Die Methode \texttt{lookup\_worker()} ist defensiv implementiert:
Bei Netzwerkfehlern oder Timeout gibt sie eine leere Liste zurück
(kein Absturz, kein Exception-Propagation):

\begin{lstlisting}[language=Python,caption={Fehlertoleranter LookupWorker-Aufruf}]
try:
    response = self._stub.LookupWorker(
        taskgrid_pb2.LookupRequest(task_type=task_type),
        timeout=5.0,
    )
except grpc.RpcError as e:
    log_event(logger, "error", "LOOKUP_WORKER_failed", error=str(e.code()))
    return []   # kein Absturz -- Dispatcher re-enqueued den Task
\end{lstlisting}

\subsubsection*{Round-Robin-Selektor}

Wenn mehrere Worker für einen Task-Typ verfügbar sind, wählt
\texttt{RoundRobinSelector} einen per Round-Robin aus (ADR-003).
Pro Task-Typ wird ein unabhängiger Zähler geführt:

\begin{lstlisting}[language=Python,caption={Round-Robin-Auswahl (worker\_selector.py)}]
with self._lock:
    idx = self._counters[task_type] % len(workers)
    self._counters[task_type] += 1
return workers[idx]
\end{lstlisting}

Bei 3 Workern ergibt sich die Sequenz: w1, w2, w3, w1, w2, w3, \ldots{}
Der Counter ist thread-sicher (Lock) und wurde durch 1000 parallele
Selektionen auf Korrektheit und gleichmäßige Verteilung getestet.

\subsection{Issue \#15/\#16: Dispatch-Loop und RESULT\_RETURN}
\label{sec:dispatch-loop}

\subsubsection*{Dispatch-Loop}

Der \texttt{DispatchLoop} läuft als Daemon-Thread (\texttt{threading.Thread})
und übernimmt den gesamten Dispatch-Ablauf:

\begin{enumerate}
  \item \texttt{dequeue(timeout=1.0)}: Blockierendes Warten auf nächsten Task.
  \item Frischen Task-Zustand aus Store lesen (könnte sich seit Einreihen geändert haben).
  \item Task an \texttt{ThreadPoolExecutor(max\_workers=10)} übergeben.
\end{enumerate}

Im Executor-Thread (\texttt{\_dispatch()}):

\begin{enumerate}
  \item \texttt{lookup\_worker()} am Namensdienst aufrufen.
  \item Falls kein Worker: \texttt{time.sleep(DISPATCH\_NO\_WORKER\_RETRY\_SECS)},
        Task re-enqueuen, Thread beenden.
  \item Zustandsübergang: QUEUED/RETRYING $\rightarrow$ DISPATCHED.
  \item Timeout-Timer starten (\texttt{threading.Timer}).
  \item \texttt{ExecuteTask} per gRPC an Worker senden (\texttt{worker\_client.py}).
  \item Bei Akzeptanz (ACK ok): DISPATCHED $\rightarrow$ PROCESSING.
        \textbf{Timer läuft weiter} bis \texttt{ReturnResult} eintrifft.
  \item Bei Ablehnung/Netzwerkfehler: Timer abbrechen,
        DISPATCHED $\rightarrow$ FAILED.
\end{enumerate}

\subsubsection*{Wichtige Design-Entscheidung: Timer-Lifetime}

Ein häufiger Fehler in der Implementierung wäre, den Timeout-Timer
bereits bei Worker-Akzeptanz abzubrechen.
Der Timer läuft absichtlich weiter:

\begin{lstlisting}[language=Python,caption={Korrekte Timer-Behandlung in \_dispatch()}]
if success:
    # Worker hat Task akzeptiert -> DISPATCHED -> PROCESSING
    # Timer laeuft weiter bis ReturnResult eintrifft (Issue #18)
    transition(task, TaskState.PROCESSING)
    self._store.update(task)
else:
    # Worker abgelehnt -> sofort FAILED, Timer abbrechen
    self._cancel_timeout(task.task_id)
    transition(task, TaskState.FAILED)
\end{lstlisting}

Dieser Fehler wurde durch die Integrations-Tests aufgedeckt
(Abschnitt~\ref{sec:bugfix-timer}).

\subsubsection*{ReturnResult (Issue \#16)}

Wenn ein Worker ein Ergebnis zurücksendet, ruft er \texttt{ReturnResult} auf dem
Dispatcher auf. Die Implementierung behandelt drei Sonderfälle:

\begin{table}[H]
\centering
\begin{tabular}{lll}
\toprule
\textbf{Sonderfall} & \textbf{Reaktion} & \textbf{Begründung} \\
\midrule
Unbekannte \texttt{task\_id}  & \texttt{Ack(success=False)} & Task existiert nicht \\
Task bereits terminal & \texttt{Ack(success=True)}, ignorieren & Idempotenz (§5) \\
Ungültiger Übergang & \texttt{Ack(success=False)}, Log-Error & Race Condition \\
\bottomrule
\end{tabular}
\caption{ReturnResult -- Sonderfallbehandlung}
\end{table}

Die Idempotenz ist besonders wichtig: Ein Worker könnte nach einem Retry
ein verspätetes Ergebnis senden. Der Task ist dann bereits \texttt{FAILED}
(Terminal). Das verspätete Ergebnis wird ignoriert.

\begin{lstlisting}[language=Python,caption={ReturnResult -- Idempotenz und Timer-Cancel}]
if is_terminal(task):
    return taskgrid_pb2.Ack(success=True, message="already terminal")

if self._dispatch_loop is not None:
    self._dispatch_loop.cancel_timeout(task_id)   # Timer sofort stoppen

new_state = TaskState.COMPLETED if request.status == "COMPLETED" \
            else TaskState.FAILED
transition(task, new_state)
task.result = request.result if request.status == "COMPLETED" \
              else request.error   # Elena: error statt error_msg
self._store.update(task)
\end{lstlisting}

\subsection{Issue \#17: GET\_RESULT -- Ergebnis abfragen}
\label{sec:get-result}

\texttt{GetResult} liefert den aktuellen Zustand und das Ergebnis eines Tasks:

\begin{itemize}
  \item Unbekannte \texttt{task\_id}: gRPC-Code \texttt{NOT\_FOUND}.
  \item Task nicht terminal: Status-Wert zurückgeben, \texttt{result=""}.
  \item Task \texttt{COMPLETED}: Status + Ergebnis zurückgeben.
  \item Task \texttt{FAILED}: Status + Fehlermeldung in \texttt{result}.
\end{itemize}

Der Client kann \texttt{GetResult} mehrfach aufrufen (Polling) bis der
Status terminal ist. Ein Beispiel-Polling-Loop:

\begin{lstlisting}[language=Python,caption={Client-seitiges Polling (Beispiel)}]
while True:
    resp = stub.GetResult(GetResultRequest(task_id=task_id))
    if resp.status in ("COMPLETED", "FAILED"):
        break
    time.sleep(1.0)
\end{lstlisting}

\subsection{Issue \#18: Timeout und Retry-Logik}
\label{sec:timeout-retry}

\subsubsection*{Funktionsweise}

Für jeden dispatchen Task startet ein \texttt{threading.Timer}.
Feuert der Timer, läuft folgender Callback ab:

\begin{lstlisting}[language=Python,caption={Timeout-Callback (\_on\_timeout)}]
def _on_timeout():
    fresh = self._store.get(task_id)
    if fresh is None or is_terminal(fresh):
        return   # verspaetetes Feuern ignorieren

    # DISPATCHED/PROCESSING -> TIMEOUT -> RETRYING
    transition(fresh, TaskState.TIMEOUT)
    transition(fresh, TaskState.RETRYING)
    fresh.retry_count += 1
    self._store.update(fresh)

    if fresh.retry_count >= self._max_retries:
        transition(fresh, TaskState.FAILED)
    else:
        self._queue.enqueue(fresh)   # erneut einplanen
\end{lstlisting}

\subsubsection*{Zustandsübergänge bei Timeout}

\begin{figure}[H]
\centering
\begin{lstlisting}[language={},frame=single,basicstyle=\ttfamily\footnotesize]
DISPATCHED/PROCESSING
       |
  [Timer feuert nach DISPATCH_TIMEOUT_SECONDS]
       |
       v
    TIMEOUT
       |
       v
   RETRYING
    /       \
   /         \
  v           v
DISPATCHED  FAILED
(retry_count  (retry_count
 < MAX)        >= MAX)
\end{lstlisting}
\caption{Zustandsübergänge bei Timeout (Issue \#18)}
\end{figure}

\subsubsection*{Konfigurierbare Parameter}

\begin{table}[H]
\centering
\begin{tabular}{lll}
\toprule
\textbf{Umgebungsvariable} & \textbf{Default} & \textbf{Bedeutung} \\
\midrule
\texttt{DISPATCH\_TIMEOUT\_SECONDS}      & 30 & Sekunden bis Timer feuert \\
\texttt{MAX\_RETRIES}                    & 3  & Max. Retries vor endgültigem FAILED \\
\texttt{DISPATCH\_NO\_WORKER\_RETRY\_SECONDS} & 2 & Wartezeit wenn kein Worker verfügbar \\
\bottomrule
\end{tabular}
\caption{Timeout-Konfiguration (Issue \#18)}
\end{table}

\subsubsection*{Behobener Bug (Integration Test)}
\label{sec:bugfix-timer}

Der Integration Test deckte auf, dass der Timer ursprünglich bei Worker-Akzeptanz
gecancelt wurde (falscher Implementierungsstand).
Außerdem fehlte \texttt{PROCESSING $\rightarrow$ TIMEOUT} als erlaubter Übergang
in der Zustandsmaschine. Beide Fehler wurden im Zuge der Integrationstests
behoben und sind in der State-Machine dokumentiert.

\subsection{Issue \#19: Strukturiertes Logging}
\label{sec:logging}

Alle Komponenten des Dispatchers nutzen das zentrale Modul
\texttt{src/common/logger.py}.

\subsubsection*{Log-Format}

\begin{lstlisting}[language={},caption={Strukturiertes Log-Format (§13)}]
[component] request_id=X task_id=Y event=Z key=value key=value
\end{lstlisting}

Pflichtfelder bei jedem Event: \texttt{request\_id}, \texttt{task\_id}, \texttt{event}.
Weitere Key-Value-Paare je nach Kontext.

\textbf{Beispiele aus dem laufenden System:}

\begin{lstlisting}[language={},caption={Log-Ausgaben Dispatcher (Normalfall)}]
[dispatcher] request_id=abc task_id=1 event=POST_TASK_accepted
             task_type=sum sender=client-1 status=QUEUED queue_size=1
[dispatcher.dispatch_loop] task_id=1 event=DISPATCH_sending
             task_type=sum worker_id=w1 status=DISPATCHED
[dispatcher.dispatch_loop] task_id=1 event=DISPATCH_accepted
             worker_id=w1 status=PROCESSING
[dispatcher] request_id=abc task_id=1 event=RESULT_RETURN_stored
             worker_id=w1 status=COMPLETED duration_ms=37
[dispatcher] request_id=abc task_id=1 event=GET_RESULT_queried
             status=COMPLETED sender=client-1
\end{lstlisting}

\textbf{Timeout-Fall:}

\begin{lstlisting}[language={},caption={Log-Ausgaben Dispatcher (Timeout-Fall)}]
[dispatcher.dispatch_loop] task_id=5 event=DISPATCH_TIMEOUT_fired
             worker_id=w1 retry_count=0
[dispatcher.dispatch_loop] task_id=5 event=TIMEOUT_retrying
             retry_count=1 max_retries=3
[dispatcher.dispatch_loop] task_id=5 event=TIMEOUT_max_retries_exceeded
             retry_count=3 max_retries=3
\end{lstlisting}

\subsubsection*{Datei-Logging}

Wenn die Umgebungsvariable \texttt{LOG\_DIR} gesetzt ist, schreibt jeder Logger
zusätzlich in eine Datei \texttt{LOG\_DIR/<component>.log} mit Timestamps.
In docker-compose.yml ist \texttt{LOG\_DIR=/app/logs} gesetzt,
das \texttt{logs}-Volume persistiert die Dateien über Container-Neustarts hinaus.

\subsubsection*{Propagation-Problem und Lösung}

Python-Logger propagieren Nachrichten standardmäßig aufwärts.
Bei verschachtelten Loggern (\texttt{dispatcher.dispatch\_loop} $\rightarrow$
\texttt{dispatcher} $\rightarrow$ root) würde jeder Eintrag mehrfach ausgegeben.
Behoben durch \texttt{logger.propagate = False} in \texttt{get\_logger()}.

\subsection{Issue \#20: GET\_STATUS -- Monitoring}
\label{sec:get-status}

\subsubsection*{HTTP-Endpoint (primär)}

Der Dispatcher startet beim Hochfahren einen \texttt{HttpStatusServer}-Thread
(Port \texttt{STATUS\_HTTP\_PORT}, default 8080).
Aufruf: \texttt{curl http://localhost:8080/status}

Die Antwort ist ein JSON-Objekt mit allen neun Pflichtfeldern aus der Aufgabenstellung:

\begin{lstlisting}[language={},caption={GET /status -- Beispielantwort}]
{
  "aktive_worker": 3,
  "unterstützte_tasktypen": ["sum", "reverse", "hash", "upper", "wait"],
  "offene_tasks": 2,
  "laufende_tasks": 1,
  "abgeschlossene_tasks": 47,
  "fehlgeschlagene_tasks": 2,
  "durchschnittliche_bearbeitungszeit_ms": 85.0,
  "anzahl_timeouts": 1,
  "anzahl_retries": 1
}
\end{lstlisting}

\subsubsection*{Datenquellen der Felder}

\begin{table}[H]
\centering
\begin{tabular}{ll}
\toprule
\textbf{Feld} & \textbf{Quelle} \\
\midrule
\texttt{aktive\_worker}          & LookupWorker für alle bekannten Task-Typen, dedupliziert \\
\texttt{unterstützte\_tasktypen} & Alle einzigartigen \texttt{task\_type}-Werte im Store \\
\texttt{offene\_tasks}           & TaskStore: Count(status=QUEUED) \\
\texttt{laufende\_tasks}         & TaskStore: Count(DISPATCHED + PROCESSING + RETRYING) \\
\texttt{abgeschlossene\_tasks}   & TaskStore: Count(status=COMPLETED) \\
\texttt{fehlgeschlagene\_tasks}  & TaskStore: Count(FAILED + TIMEOUT) \\
\texttt{durchschnittliche\_bearbeitungszeit\_ms} & Ø (timestamp\_completed - timestamp\_dispatched) × 1000 \\
\texttt{anzahl\_timeouts}        & Summe aller \texttt{task.retry\_count} \\
\texttt{anzahl\_retries}         & Identisch (jeder Retry durch Timeout ausgelöst) \\
\bottomrule
\end{tabular}
\caption{Monitoring-Felder und ihre Datenquellen (Issue \#20)}
\end{table}

\subsubsection*{gRPC GetStatus (sekundär)}

Zusätzlich zum HTTP-Endpoint ist \texttt{GetStatus} in der gRPC-Schnittstelle
implementiert. Er gibt \texttt{queued\_tasks}, \texttt{active\_workers}
direkt als Proto-Felder zurück sowie das vollständige JSON als
\texttt{details}-String -- ohne Proto-Schemaänderung.

\subsubsection*{Designentscheidung (ADR-006)}

HTTP wurde gegenüber einer Proto-Erweiterung bevorzugt,
da §8 der Aufgabenstellung keine Proto-Änderungen für Monitoring erzwingt,
HTTP universell mit \texttt{curl} oder Browser abrufbar ist,
und die Elena-Proto-Kompatibilität damit vollständig erhalten bleibt
(kein Regenerierungsrisiko).

\subsection{Issue \#21: Queue-Verwaltung und Nebenläufigkeit}
\label{sec:concurrency}

\subsubsection*{Thread-Modell des Dispatchers}

Im Betrieb laufen mindestens vier Thread-Typen parallel:

\begin{table}[H]
\centering
\begin{tabular}{lll}
\toprule
\textbf{Thread-Typ} & \textbf{Anzahl} & \textbf{Aufgabe} \\
\midrule
gRPC-Server-Threads    & bis 10 & PostTask, ReturnResult, GetResult bedienen \\
DispatchLoop-Thread    & 1      & Queue lesen, Executor befüllen \\
Executor-Threads       & bis 10 & ExecuteTask an Worker senden \\
Timeout-Timer-Threads  & 1 pro task & Nach DISPATCH\_TIMEOUT\_SECONDS feuern \\
HTTP-Server-Thread     & 1      & GET /status bedienen \\
\bottomrule
\end{tabular}
\caption{Thread-Modell des Dispatchers}
\end{table}

\subsubsection*{Locking-Strategie (ADR-007)}

Jede geteilte Datenstruktur schützt sich selbst:

\begin{itemize}
  \item \textbf{TaskStore}: \texttt{threading.Lock()} auf allen Methoden.
        \texttt{all\_tasks()} und \texttt{worker\_load\_map()} erzeugen unter Lock
        einen Snapshot -- der Aufrufer erhält einen konsistenten Zustand.
  \item \textbf{TaskQueue}: \texttt{queue.Queue} aus der Python-Standardbibliothek
        (intern thread-sicher, blockierendes \texttt{get(timeout=...)}).
  \item \textbf{RoundRobinSelector}: \texttt{threading.Lock()} auf dem
        Read-Modify-Write-Zähler.
  \item \textbf{new\_task\_id()}: \texttt{threading.Lock()} auf globalem Zähler.
  \item \textbf{\_timers-Dict}: \texttt{\_timers\_lock} verhindert Race zwischen
        Timer-Start (Executor-Thread) und Timer-Cancel (gRPC-Thread).
\end{itemize}

\subsubsection*{Worker-Last-Tracking}

\texttt{TaskStore.worker\_load\_map()} liefert einen Snapshot der aktuellen
Auslastung pro Worker -- gezählt werden Tasks in DISPATCHED oder PROCESSING
mit gesetztem \texttt{assigned\_worker}:

\begin{lstlisting}[language=Python,caption={worker\_load\_map() -- aktive Tasks pro Worker}]
active = (TaskState.DISPATCHED, TaskState.PROCESSING)
with self._lock:
    load = {}
    for t in self._tasks.values():
        if t.status in active and t.assigned_worker:
            load[t.assigned_worker] = load.get(t.assigned_worker, 0) + 1
    return load
\end{lstlisting}

Diese Methode wird vom \texttt{StatusCollector} für das Monitoring genutzt
und im Integrations-Test auf Konsistenz unter Last geprüft.

\subsubsection*{Nachgewiesene Thread-Sicherheit}

Die Nebenläufigkeit wurde durch automatisierte Stress-Tests in
\texttt{tests/test\_concurrency.py} nachgewiesen:

\begin{table}[H]
\centering
\begin{tabular}{llr}
\toprule
\textbf{Test} & \textbf{Was wird geprüft} & \textbf{Threads} \\
\midrule
Concurrent enqueue/dequeue & Kein Task verloren (1000 Tasks) & 20 \\
Concurrent add (Store) & Keine Duplikate (500 Tasks) & 10 \\
Concurrent read/write (Store) & Keine Exception & 8 \\
Round-Robin unter Last & Gleichmäßige Verteilung (1000 Selects) & 10 \\
Concurrent PostTask (live) & Alle 100 Tasks gespeichert & 100 \\
Concurrent ReturnResult (live) & Nur erster Wert gilt (Idempotenz) & 20 \\
\bottomrule
\end{tabular}
\caption{Nebenläufigkeits-Tests (Issue \#21)}
\end{table}

\subsection{ADR-004: Speicherung von Tasks und Ergebnissen}
\label{sec:adr004}

\textbf{Status}: Akzeptiert \quad
\textbf{Betrifft}: \texttt{src/dispatcher/task\_store.py}

\subsubsection*{Kontext}

Der Dispatcher benötigt einen Speicher für Tasks während ihrer Lebensdauer.
Die Aufgabenstellung (§8) verbietet explizit, gemeinsame Datenbanken oder
Dateisysteme für die Task-Übergabe zu verwenden.

\subsubsection*{Entscheidung}

\textbf{In-Memory-Speicherung}: Alle Tasks in einem Python-\texttt{dict} im RAM.
Kein externer Dienst, kein Docker-Volume für Task-Daten.

\subsubsection*{Bewertete Alternativen}

\begin{table}[H]
\centering
\begin{tabular}{lp{4.5cm}p{4.5cm}}
\toprule
\textbf{Alternative} & \textbf{Argument dagegen} & \textbf{Folge} \\
\midrule
SQLite/PostgreSQL &
  §8 verbietet Shared-DB; externer Container; Latenz &
  Abgelehnt \\
Redis &
  Weiterer Container; Serialisierung nötig; horizontale Skalierung des Dispatchers nicht gefordert &
  Abgelehnt \\
JSON-Datei &
  §8: Volumes nur für Logs; atomares Schreiben komplex; doppeltes Locking &
  Abgelehnt \\
Hybrid (RAM + Snapshot) &
  Kein klarer Anforderungs-Vorteil; Durability nicht gefordert &
  Abgelehnt \\
\textbf{In-Memory} (gewählt) &
  -- &
  Akzeptiert \\
\bottomrule
\end{tabular}
\caption{ADR-004: Bewertete Alternativen für Task-Speicherung}
\end{table}

\subsubsection*{Konsequenzen}

Positiv: O(1)-Zugriff, keine Abhängigkeiten, konform mit §8,
vollständig mock-testbar (95 Unit-Tests).

Negativ/Akzeptiert: Tasks gehen bei Dispatcher-Neustart verloren.
Dies ist für den Studienkontext akzeptabel -- Durability ist keine Anforderung.

\subsection{Konfigurationsreferenz}

Alle Dispatcher-Parameter werden ausschließlich über Umgebungsvariablen
konfiguriert. Dies ermöglicht die Anpassung ohne Code-Änderung:

\begin{table}[H]
\centering
\begin{tabular}{lll}
\toprule
\textbf{Variable} & \textbf{Default} & \textbf{Bedeutung} \\
\midrule
\texttt{DISPATCHER\_PORT}                    & 50051             & gRPC-Port \\
\texttt{STATUS\_HTTP\_PORT}                  & 8080              & HTTP Monitoring-Port \\
\texttt{NAMENSDIENST\_HOST}                  & \texttt{namensdienst} & Namensdienst-Hostname \\
\texttt{NAMENSDIENST\_PORT}                  & 50052             & Namensdienst-Port \\
\texttt{DISPATCH\_TIMEOUT\_SECONDS}          & 30                & Sekunden bis Timeout \\
\texttt{MAX\_RETRIES}                        & 3                 & Max. Retries vor FAILED \\
\texttt{DISPATCH\_NO\_WORKER\_RETRY\_SECONDS} & 2                & Wartezeit kein Worker \\
\texttt{LOG\_DIR}                            & \texttt{logs}     & Verzeichnis für Log-Dateien \\
\bottomrule
\end{tabular}
\caption{Dispatcher -- Konfigurationsreferenz}
\end{table}

\subsection{Testergebnisse}

\subsubsection*{Unit-Tests}

\begin{table}[H]
\centering
\begin{tabular}{llr}
\toprule
\textbf{Testdatei} & \textbf{Abdeckung} & \textbf{Tests} \\
\midrule
\texttt{test\_post\_task.py}       & Issue \#13: Validierung, Speicherung, Queue & 8 \\
\texttt{test\_get\_result.py}      & Issue \#17: Alle Statuszustände, NOT\_FOUND & 8 \\
\texttt{test\_receive\_result.py}  & Issue \#16: COMPLETED, FAILED, Idempotenz & 9 \\
\texttt{test\_dispatch.py}         & Issue \#15: Dispatch-Flow, Timer, Retry & 13 \\
\texttt{test\_state\_machine.py}   & Alle 8 Zustände, Übergänge, Fehlerfälle & 14 \\
\texttt{test\_logging.py}          & Issue \#19: Format, Datei-Logging & 13 \\
\texttt{test\_get\_status.py}      & Issue \#20: Alle 9 Felder, gRPC + HTTP & 11 \\
\texttt{test\_concurrency.py}      & Issue \#21: Race Conditions, Stress-Tests & 13 \\
\midrule
\textbf{Gesamt}                    &                                           & \textbf{89} \\
\bottomrule
\end{tabular}
\caption{Dispatcher -- Unit-Test-Übersicht (89/89 bestanden)}
\end{table}

\subsubsection*{Integrations-Tests}

Das Skript \texttt{integration\_test.py} startet einen echten gRPC-Server
und HTTP-Server (ohne Mocks) und führt 49 End-to-End-Tests durch:
alle Issues \#13--\#21, Timeout-Szenarien mit echtem Timer,
vollständiger Task-Durchlauf (QUEUED $\rightarrow$ COMPLETED).

\textbf{Ergebnis}: 49/49 Tests bestanden.

Dabei wurden zwei Implementierungsfehler entdeckt und behoben:
\begin{enumerate}
  \item Timer wurde bei Worker-Akzeptanz gecancelt statt erst bei \texttt{ReturnResult}.
  \item \texttt{PROCESSING $\rightarrow$ TIMEOUT} fehlte in der Zustandsmaschine.
\end{enumerate}
